\import{surreal/real-approximation.tex}

\section{Dyadic intervals for surreal reals}

\begin{lemma}\label{successor_omega_surreal_symmetric_natural_bound}
    Suppose $x\in\surrealStage{\suc{\naturals}}$.
    Suppose
        $\surrealLt{\surrealNeg{\ordinalSurreal{\naturals}}}{x}$ and
        $\surrealLt{x}{\ordinalSurreal{\naturals}}$.
    Then there exists a natural number $N$ such that
        $\surrealLt{\surrealNeg{\ordinalSurreal{N}}}{x}$ and
        $\surrealLt{x}{\ordinalSurreal{N}}$.
\end{lemma}
\begin{proof}
    Take a natural number $A$ such that
        $\surrealLt{\surrealNeg{\ordinalSurreal{A}}}{x}$
        by \obv{successor_omega_surreal_natural_bound_below}.
    Take a natural number $B$ such that
        $\surrealLt{x}{\ordinalSurreal{B}}$
        by \obv{successor_omega_surreal_natural_bound_above}.
    $A\in B$ or $A=B$ or $B\in A$
        by
        \obv{ordinal_in_trichotomy,natural_number_is_ordinal}.
    \begin{byCase}
        \caseOf{$A=B$.}
            Follows by \auto.
        \caseOf{$A\neq B$.}
            \begin{byCase}
                \caseOf{$A\in B$.}
                    $\surrealLt{\ordinalSurreal{A}}
                        {\ordinalSurreal{B}}$
                        by
                        \obv{ordinal_surreal_lt_iff,natural_number_is_ordinal}.
                    Hence
                        $\surrealLt
                            {\surrealNeg{\ordinalSurreal{B}}}
                            {\surrealNeg{\ordinalSurreal{A}}}$
                        by
                        \obv{surreal_negation_lt_reverses,ordinal_surreal_is_surreal,natural_number_is_ordinal}.
                    Hence
                        $\surrealLt
                            {\surrealNeg{\ordinalSurreal{B}}}{x}$
                        by
                        \auto{surreal_lt_transitive,surreal_negation_is_surreal,ordinal_surreal_is_surreal,natural_number_is_ordinal,surreal_stage_characterization,suc_ordinal,omega_is_ordinal}.
                \caseOf{$A\notin B$.}
                    Hence $B\in A$.
                    $\surrealLt{\ordinalSurreal{B}}
                        {\ordinalSurreal{A}}$
                        by
                        \obv{ordinal_surreal_lt_iff,natural_number_is_ordinal}.
                    Hence $\surrealLt{x}{\ordinalSurreal{A}}$
                        by
                        \auto{surreal_lt_transitive,ordinal_surreal_is_surreal,natural_number_is_ordinal,surreal_stage_characterization,suc_ordinal,omega_is_ordinal}.
            \end{byCase}
    \end{byCase}
\end{proof}

\begin{theorem}\label{surreal_real_dyadic_interval}
    Suppose $x\in\reals$.
    Let $n$ be an ordinal.
    Suppose $n$ is a natural number.
    Then there exists $q\in\surrealStage{\naturals}$ such that
        $\surrealLt{q}{x}$ and
        $\surrealLt{x}{\surrealAdd{q}{\surrealEpsilon{n}}}$.
\end{theorem}
\begin{proof}[Proof by transfinite induction]
    \begin{byCase}
        \caseOf{$x\in\surrealStage{\naturals}$.}
            Follows by \obv{finite_surreal_dyadic_interval}.
        \caseOf{$x\notin\surrealStage{\naturals}$.}
            $x\in\surrealStage{\suc{\naturals}}$ and
                $\surrealLt
                    {\surrealNeg{\ordinalSurreal{\naturals}}}{x}$ and
                $\surrealLt{x}{\ordinalSurreal{\naturals}}$
                by \obv{surreal_real_elim,surreal_real_bounds}.
            $n=\emptyset$ or there exists a natural number $k$ such that
                $n=\suc{k}$ by \obv{natural_zero_or_successor}.
            \begin{byCase}
                \caseOf{$n=\emptyset$.}
                    Take a natural number $N$ such that
                        $\surrealLt
                            {\surrealNeg{\ordinalSurreal{N}}}{x}$ and
                        $\surrealLt{x}{\ordinalSurreal{N}}$
                        by
                        \obv{successor_omega_surreal_symmetric_natural_bound}.
                    Take $q\in\surrealStage{\naturals}$ such that
                        $\surrealLt{q}{x}$ and
                        $\surrealLt{x}{\surrealAdd{q}{\surrealOne}}$
                        by
                        \obv{nonfinite_surreal_unit_interval_in_symmetric_bound,natural_number_is_ordinal}.
                    $\surrealEpsilon{n}=\surrealOne$
                        by \obv{surreal_epsilon_zero}.
                    Follows by \auto.
                \caseOf{$n\neq\emptyset$.}
                    Take $k\in\naturals$ such that $n=\suc{k}$.
                    $n=\suc{k}$.
                    $k\in n$ by \obv{suc_intro_self}.
                    Take $q\in\surrealStage{\naturals}$ such that
                        $\surrealLt{q}{x}$ and
                        $\surrealLt{x}
                            {\surrealAdd{q}{\surrealEpsilon{k}}}$.
                    Let $m=\surrealAdd{q}{\surrealEpsilon{n}}$.
                    $\surrealEpsilon{n}\in
                        \surrealStage{\naturals}$
                        by \obv{surreal_epsilon_in_finite_stage}.
                    Hence $m\in\surrealStage{\naturals}$
                        by \auto{surreal_addition_finite_stage}.
                    $\surreal{x}$ and $\surreal{m}$
                        by
                        \auto{surreal_real_is_surreal,surreal_stage_characterization,omega_is_ordinal}.
                    $\surrealLt{x}{m}$ or $x=m$ or
                        $\surrealLt{m}{x}$
                        by \obv{surreal_lt_trichotomy}.
                    \begin{byCase}
                        \caseOf{$x=m$.}
                            Contradiction.
                        \caseOf{$x\neq m$.}
                            \begin{byCase}
                                \caseOf{$\surrealLt{x}{m}$.}
                                    Follows by \auto.
                                \caseOf{not $\surrealLt{x}{m}$.}
                                    Hence $\surrealLt{m}{x}$.
                                    $\surrealAdd{\surrealEpsilon{n}}
                                        {\surrealEpsilon{n}}
                                        =\surrealEpsilon{k}$
                                        by
                                        \auto{surreal_epsilon_halving,natural_number_is_ordinal}.
                                    Hence
                                        $\surrealAdd{m}
                                            {\surrealEpsilon{n}}
                                        =\surrealAdd{q}
                                            {\surrealEpsilon{k}}$
                                        by
                                        \auto{surreal_add_associative,surreal_add_is_surreal,surreal_epsilon_is_surreal,surreal_stage_characterization,omega_is_ordinal}.
                            \end{byCase}
                    \end{byCase}
            \end{byCase}
    \end{byCase}
\end{proof}
